\def\xhref#1#2{\href{#1}{\hskip16pt}\llap{#2}}
\def\xxhref#1#2{\href{#1}{\hskip14pt}\llap{#2}}
\makeatletter
\def\@oddhead{\hfill}
\def\@evenhead{\hfill}
\makeatother
   
   
\normalsize
%\fontsize{9.6}{11.6}\selectfont
%\parskip4pt plus2pt minus1pt

%\advance\baselineskip by.1pt

\pznspis{Znajdując błyszczący pierścień w~stogu siana, ups...\ w~kosmosie}


\pzntyt{Znajdując błyszczący pierścień w~stogu siana, ups\dots\ w~kosmosie}

\vspace*{4pt}
\begin{dwieszpalty}\vspace*{-10pt}
\marg{
    \includegraphics[clip,trim=90 50 100 70,width=0.95\linewidth]{potm2406a_not_inverted_small.jpg}\\[4pt]
\textcolor{black}{Rys. 1.    {Masywna galaktyka (pomarańczowa kropka w~centrum pierścienia) zniekształca światło emitowane przez kwazara (błyszczące diamenty) w~zjawisku zwanym soczewkowaniem grawitacyjnym. Źródło zdjęcia: NASA/ESA/CSA}}%
}Powyższe zdjęcie, do złudzenia przypominające piękny pierścień, zostało wykonane przez Kosmiczny Teleskop Jamesa Webba prawie rok temu, jednak wciąż zachwyca! Ponadto dostarcza ważne informacje o~czarnych dziurach.

Cztery błyszczące punkty, przypominające drogocenne kamienie pierścienia, to zwielokrotniony obraz znanego już wcześniej astronomom kwazara o~mało romantycznej nazwie RX~J1131-1231. Znajduje się on około 6 miliardów lat świetlnych od Ziemi i~był wcześniej obserwowany przez rentgenowski teleskop Chandra i~Kosmiczny Teleskop Hubble'a (rys. 2). Jednak nigdy wcześniej nie widzieliśmy go z~taką dokładnością i~nie błyszczał tak pięknie.
\marg{Kwazarem nazywamy niezwykle jasne aktywne jądro galaktyczne. Pomimo niewielkich rozmiarów, porównywalnych do rozmiaru naszego Układu Słonecznego, kwazary są w~stanie oświetlić całe galaktyki i~przyćmić światło gwiazd. Energia, którą produkują i~wyświecają, pochodzi z~obszaru znajdującego się wokół supermasywnej czarnej dziury. } 

Jasny obiekt w~środku pierścienia to galaktyka eliptyczna, której masa (wraz z~otaczającą ją ciemną materią) działa jak soczewka. Dzięki niej światło zmierzające w~naszą stronę od kwazara, znajdującego się za naszą masywną eliptyczną galaktyką, zostało zakrzywione i~rozszczepione. Zjawisko~to, nazywane soczewkowaniem grawitacyjnym, po raz pierwszy zostało przewidziane przez Einsteina. Oferuje ono rzadką możliwość badania obiektów astronomicznych, działając jak naturalny teleskop, powielając i~powiększając obraz obiektu, którego często nie bylibyśmy w~stanie zaobserwować. W~tym wypadku zwielokrotniony obraz pozwolił naukowcom zbadać obszar znajdujący się wokół czarnej dziury we wnętrzu kwazara.
\marg{\includegraphics[clip,trim=0 10 0 20,width=.95\linewidth]{Space_telescopes_measure_black_hole's_spin_cropped.jpg}\\[4pt]
    \textcolor{black}{Rys. 2. Zwielokrotniony obraz kwazara widoczny w~danych pochodzących z~rentgenowskiego teleskopu Chandra (różowy kolor).  Dane optyczne Hubble'a w~kolorach czerwonym, zielonym i~niebieskim pokazują galaktykę eliptyczną w~środku obrazu. Żródło: X-ray: NASA/CXC/Univ of Michigan/R.C.Reis et al; Optical: NASA/STScl}}
\end{dwieszpalty}

Przedstawiciele Europejskiej Agencji Kosmicznej (European Space Agency, ESA) w~oświadczeniu dołączonym do opublikowanego zdjęcia wyjaśniają, że pomiary emisji promieniowania rentgenowskiego z~kwazarów pozwalają określić, jak szybko obraca się centralna czarna dziura w~jego wnętrzu. To z~kolei daje naukowcom ważne wskazówki na temat tego, jak takie czarne dziury powstają, rosną i~przybierają na masie wraz z~upływem czasu. 
Dane obserwacyjne uzyskane z~Kosmicznego Teleskopu Jamesa Webba, a~także wcześniejsze obserwacje wykonane za pomocą teleskopów Chandra i~Hubble'a umożliwiły bardzo dokładną analizę kwazara i~jednocześnie  analizę tempa wzrostu znajdującej się w~nim czarnej dziury. 
Okazało się, że czarna dziura w~tym konkretnym kwazarze obraca się z~prędkością ponad połowy prędkości światła! Z~tego można wnioskować, że urosła w~wyniku fuzji, czyli zderzenia się z~inną czarną dziurą, a~nie tylko w~wyniku przyciągania i~pochłaniania materii z~różnych kierunków.

RX~J1131-1231 jest uważany za jeden z~najlepiej przesoczewkowanych kwazarów odkrytych do tej pory. Jest też najbardziej błyszczącym z~kosmicznych pierścieni!





\rightline{\large
%\textit{Anna DURKALEC}%
%\textit{Michał BEJGER}%
\textit{Katarzyna MAŁEK}%
}%
%}

\medskip

%{\scriptsize \rightline{Centrum Astronomiczne im. Mikołaja Kopernika PAN,}
%   \rightline  {Istituto Nazionale di Fisica Nucleare (INFN), Sezione di Ferrara, Włochy}\par} 

%{\scriptsize \rightline{Zakład Astrofizyki, Departament Badań Podstawowych, Narodowe Centrum Badań Jądrowych}}  %A.Durkalec

 \rightline{ \scriptsize  Departament Badań Podstawowych (BP4), Zakład Astrofizyki, Narodowe Centrum Badań Jądrowych}%K.Małek
%\vfill
% \end{dwieszpalty}


%\normalsize
%\input 13-niebo

\eject

\endinput
\begin{figure*}[ht!]
    \centering
    \includegraphics[width=0.6\linewidth]{potm2406a_not_inverted_small.jpg}
    \caption{ Zdjęcia bez odwrócenia kolorów}
\end{figure*}